% !TeX root = ../数学分析习题（裴礼文）笔记.tex
% 由主文档通过 \input 载入；请编译主文档。
\setcounter{chapter}{3} % 保留原书章号。
\chapter{一元函数积分学}
\label{chap:pei-int-integral-calculus}

本章对应教材第四章，围绕积分与极限的互相转化、Riemann 可积性的判别、定积分估计与中值定理、经典积分不等式，以及反常积分的计算和收敛性展开。核心方法是先辨认可用的结构，再用积分和、振幅、单调性或尾部估计控制误差。有限区间上的积分默认指 Riemann 积分；在出现无界被积函数或无穷区间时，另按反常积分定义处理。除另有说明，区间均非退化、函数均取实值；有界性、非负性、连续性与绝对收敛性分别在所需位置注明，不能相互替代。

% 教材来源：数学分析/textbooks/数分习题（分章原文）/数分习题-4.pdf。
% 覆盖五节的主要结论、证明方法与代表性例题；按主题整合，非逐题转录。

\section{积分与极限}
\label{sec:pei-int-integral-limits}
% 教材对应：§4.1，PDF 第 1--17 页（原书第 250--266 页）。

\subsection{把和式、乘积与误差还原为积分}
\begin{theorem}
\label{thm:pei-int-riemann-sum}
设 $f$ 在 $[a,b]$ 上 Riemann 可积。对分划 $P_n:a=x_{n,0}<\cdots<x_{n,m_n}=b$，记网格宽度 $\lVert P_n\rVert=\max_i(x_{n,i}-x_{n,i-1})$。若 $\lVert P_n\rVert\to0$，则不论标记点 $\xi_{n,i}\in[x_{n,i-1},x_{n,i}]$ 如何选取，都有
\[
 \sum_{i=1}^{m_n}f(\xi_{n,i})(x_{n,i}-x_{n,i-1})\longrightarrow\int_a^b f(x)\,dx.
\]
\end{theorem}
识别积分和时，先确定区间与每项中的小区间长度，再确定函数值对应的标记点；\emph[caution]{分划不必等分，但最大的小区间长度必须趋于零}。例如 $\alpha>0$ 时，$n^{-\alpha-1}\sum_{k=1}^n k^\alpha\to1/(\alpha+1)$；若 $-1<\alpha<0$，被积函数 $x^\alpha$ 在零点无界，不能直接援引上述定理，须先截去零点附近的一段，再用单调性估计该段和式。这样仍得到同一极限。遇到乘积时先取对数，把乘积变成和式；取对数要求各因子为正，最后还要由指数函数的连续性还原极限。

\begin{example}
\label{ex:pei-int-geometric-partition}
设 $b\ge1$，求 $\displaystyle\lim_{n\to\infty}(b^{1/n}-1)\sum_{k=0}^{n-1}b^{k/n}\sin(b^{k/n})$。
\end{example}
\begin{solve}
当 $b=1$ 时原式恒为零。当 $b>1$ 时令 $x_k=b^{k/n}$，则 $x_{k+1}-x_k=(b^{1/n}-1)b^{k/n}$，而网格宽度不超过 $b(b^{1/n}-1)\to0$。原式为 $\sin x$ 在 $[1,b]$ 上的左端点积分和，故极限为 $\int_1^b\sin x\,dx=\cos1-\cos b$。
\end{solve}

\begin{proposition}
\label{prop:pei-int-quadrature-asymptotic}
设 $h=(b-a)/n$，$R_n=h\sum_{i=1}^n f(a+ih)$，$M_n=h\sum_{i=1}^n f(a+(i-\tfrac12)h)$。若 $f\in C^1[a,b]$，则
\begin{equation}
 \lim_{n\to\infty}n\left(R_n-\int_a^b f\right)=\frac{b-a}{2}\bigl(f(b)-f(a)\bigr).
 \label{eq:pei-int-right-sum-error}
\end{equation}
若 $f\in C^2[a,b]$，则
\begin{equation}
 \lim_{n\to\infty}n^2\left(\int_a^b f-M_n\right)=\frac{(b-a)^2}{24}\bigl(f'(b)-f'(a)\bigr).
 \label{eq:pei-int-midpoint-error}
\end{equation}
\end{proposition}
\begin{proof}
在每个右端点 $x_i$ 附近，由 $f'$ 的一致连续性，$f(x)=f(x_i)+f'(x_i)(x-x_i)+o(h)$，余项关于区间编号一致。逐段积分并求和，得到 $R_n-\int_a^b f=\tfrac{h^2}{2}\sum_i f'(x_i)+o(h)$，再应用积分和极限与 Newton--Leibniz 公式，即得式~\eqref{eq:pei-int-right-sum-error}。在中点 $m_i$ 展开至二阶，奇次项的积分为零，且 $\int_{x_{i-1}}^{x_i}(x-m_i)^2\,dx=h^3/12$，故 $\int_a^b f-M_n=\tfrac{h^3}{24}\sum_i f''(m_i)+o(h^2)$，由此得式~\eqref{eq:pei-int-midpoint-error}。余项的一致性保证 $n$ 个小区间的误差相加后仍具有所需阶数。
\end{proof}
例如取 $f(x)=1/(1+x)$、$[a,b]=[0,1]$，则 $n(\ln2-R_n)\to1/4$，$n^2(\ln2-M_n)\to1/32$。这类极限的目标是误差的首项，单独证明积分和收敛还不够。

\subsection{先分区间，再交换极限与积分}
若 $f_n\to f$ 在有限区间上一致成立，则 $|\int f_n-\int f|\le(b-a)\sup|f_n-f|\to0$。只有逐点收敛时，这个结论不自动成立；例如 $f_n(x)=n^2x(1-x)^n$ 在 $[0,1]$ 上逐点趋于零，但其积分为 $n^2/((n+1)(n+2))\to1$。教材中的许多例题采用\emph[tip]{小区间上控制长度，大区间上取得一致估计}，从而在 Riemann 积分框架内直接证明极限。

\begin{example}
证明 $\displaystyle\int_0^{\pi/2}\sin^n x\,dx\to0$。
\end{example}
\begin{solve}
固定 $0<\delta<\pi/2$，有
\[
 0\le\int_0^{\pi/2}\sin^n x\,dx
 \le\frac\pi2(\cos\delta)^n+\delta.
\]
先令 $n\to\infty$，得到上极限不超过 $\delta$，再令 $\delta\downarrow0$。这里先选并固定 $\delta$，再选 $n$，避免把依赖于 $n$ 的估计误当成一致估计。
\end{solve}

\begin{theorem}
\label{thm:pei-int-concentrating-kernel}
设 $K_n\ge0$ 在 $[a,b]$ 上可积，$\int_a^bK_n\to A<\infty$，且对每个 $\delta>0$，$\int_{[a,b]\setminus(x_0-\delta,x_0+\delta)}K_n\to0$。若 $f$ 在 $[a,b]$ 上有界、Riemann 可积，并在 $x_0\in[a,b]$ 连续，则
\[
 \int_a^b f(x)K_n(x)\,dx\longrightarrow Af(x_0).
\]
\end{theorem}
\begin{proof}
写成 $\int fK_n-f(x_0)\int K_n=\int(f-f(x_0))K_n$。在 $x_0$ 的 $\delta$ 邻域内利用连续性把 $|f-f(x_0)|$ 控制在任意给定的小量之下；在其余部分利用 $f$ 的有界性及核的质量趋于零。由于 $\int K_n$ 有界，两部分的误差都能趋于零；再用总质量的极限即得结论。
\end{proof}
这一结论称为集中核的逼近性质，关键在核的非负性、总质量和远离集中点后的质量估计。它统一了教材中“先减去待求常数，再估计误差”的拟合法。例如对连续的 $f$，
\begin{align*}
 \lim_{h\downarrow0}\int_0^1\frac{h f(x)}{h^2+x^2}\,dx&=\frac\pi2f(0),\\
 \lim_{n\to\infty}(n+1)\int_0^1x^nf(x)\,dx&=f(1),\\
 \lim_{n\to\infty}\frac{n+1}{2}\int_{-1}^1(1-|x|)^nf(x)\,dx&=f(0).
\end{align*}
第一式的核总质量为 $\arctan(1/h)\to\pi/2$，在 $[\delta,1]$ 上的积分不超过 $h/\delta$；第二式在 $[0,1-\delta]$ 上的质量为 $(1-\delta)^{n+1}$；第三式的总质量恒为一，远离零点后也具有几何衰减。各式因而无需预先假设 $f$ 可微。

\subsection{高次幂把积分集中到最大值附近}
\begin{theorem}
\label{thm:pei-int-power-integral}
设 $f\ge0$、$g>0$ 在 $[a,b]$ 上连续，$a<b$。则
\begin{equation}
 \lim_{n\to\infty}\left(\int_a^bf(x)^ng(x)\,dx\right)^{1/n}
 =\max_{[a,b]}f.
 \label{eq:pei-int-lp-max}
\end{equation}
若进一步设 $f$ 严格递减、$f(a)=1$、$f(b)=0$，记 $I_n=\int_a^bf(x)^n\,dx$，则 $I_{n+1}/I_n\to1$。
\end{theorem}
\begin{proof}
令 $M=\max f$。$M=0$ 时第一式成立；$M>0$ 时，上界为 $M(\int g)^{1/n}\to M$。对任意 $0<\varepsilon<M$，最大值点的某个单侧或双侧邻域 $J$ 具有正长度，且 $f>M-\varepsilon$；因 $\int_Jg>0$，下界为 $(M-\varepsilon)(\int_Jg)^{1/n}\to M-\varepsilon$。令 $\varepsilon\downarrow0$ 即得第一式。

对比值结论，固定 $a<c<d<b$，有
\[
 \frac{\int_d^bf^n}{I_n}\le\frac{b-d}{c-a}\left(\frac{f(d)}{f(c)}\right)^n\longrightarrow0.
\]
故 $\liminf I_{n+1}/I_n\ge f(d)$，而上界为一。令 $d\downarrow a$，便得到比值极限。这个证明还说明幂权重在最大值点附近集中；比值极限并非仅由 $I_n^{1/n}\to1$ 推出。
\end{proof}
若 $f\ge0$ 连续且严格递增，积分中值点 $x_n$ 由 $f(x_n)^n=(b-a)^{-1}\int_a^bf(x)^n\,dx$ 确定，则式~\eqref{eq:pei-int-lp-max} 给出 $f(x_n)\to f(b)$。严格单调性与连续性遂推出 $x_n\to b$；中值点的位置需要结合函数的单调性判断。

\subsection{快速振荡的平均作用与平移连续性}
\begin{theorem}
\label{thm:pei-int-periodic-average}
设 $f$ 在 $[a,b]$ 上 Riemann 可积，$g$ 以 $T>0$ 为周期且在 $[0,T]$ 上 Riemann 可积，记 $\bar g=T^{-1}\int_0^Tg(t)\,dt$。则
\begin{equation}
 \lim_{\lambda\to+\infty}\int_a^bf(x)g(\lambda x)\,dx
 =\bar g\int_a^bf(x)\,dx.
 \label{eq:pei-int-periodic-average}
\end{equation}
特别地，$\int_a^bf(x)\sin(\lambda x)\,dx\to0$，余弦的结论相同，这就是此处的 Riemann--Lebesgue 引理。
\end{theorem}
\begin{proof}
先设 $f$ 为一个区间的示性函数。将积分区间分成完整周期与至多两个残段，完整周期的积分等于区间长度乘以 $\bar g$，残段的误差为 $O(\lambda^{-1})$；由线性性，阶梯函数也成立。一般可积函数可用阶梯函数 $s$ 在积分意义下逼近，使 $\int_a^b|f-s|$ 任意小，而
\[
 \left|\int_a^b(f-s)g(\lambda x)\,dx\right|
 \le\sup|g|\int_a^b|f-s|\,dx.
\]
右边与 $\lambda$ 无关，故可先固定 $s$，再取极限，最后令逼近误差趋于零。$g$ 允许变号；也可将其分解为正部与负部分别处理。
\end{proof}

\begin{proposition}
\label{prop:pei-int-translation}
将 $[a,b]$ 上的 Riemann 可积函数 $f$ 在区间外延拓为零，则
\[
 \lim_{h\to0}\int_{\mathbb R}|f(x+h)-f(x)|\,dx=0.
\]
若 $f$ 在较大的区间上连续，则还有
\[
 \lim_{h\to0}\frac1h\int_a^b\bigl(f(x+h)-f(x)\bigr)\,dx=f(b)-f(a).
\]
\end{proposition}
\begin{proof}
第一式先对阶梯函数证明：平移后只在有限个分点附近发生变化，故积分为 $O(|h|)$。选阶梯函数 $s$ 使 $\int|f-s|<\varepsilon$，再由平移换元与三角不等式，原积分不超过 $2\varepsilon+\int|s(x+h)-s(x)|$，于是得证。第二式把分子改写为 $\int_b^{b+h}f-\int_a^{a+h}f$，除以 $h$ 后由端点连续性取极限即可，并不要求 $f'$ 存在。
\end{proof}

\section{定积分的可积性}
\label{sec:pei-int-integrability}
% 教材对应：§4.2，PDF 第 17--32 页（原书第 266--281 页）。

\subsection{定义中的量词与 Darboux 判别}
\begin{definition}
设 $f:[a,b]\to\mathbb R$，$a<b$。若存在实数 $I$，使对每个 $\varepsilon>0$，都存在 $\delta>0$，对任意满足 $\lVert P\rVert<\delta$ 的分划及任意标记点均有 $|\sum_i f(\xi_i)\Delta x_i-I|<\varepsilon$，则称 $f$ \emph{Riemann 可积（Riemann integrable）}，并记 $I=\int_a^bf$。这里 $\delta$ 只依赖于 $\varepsilon$ 与 $f$，不能依赖于其后任选的分划或标记点。
\end{definition}
可积函数必有界：若不然，任取一个足够细的固定分划，至少一个小区间上函数无界；固定其余标记点而改变这个小区间的标记点，就能使积分和无界，与定义矛盾。以下讨论均设 $f$ 有界。对 $P:a=x_0<\cdots<x_m=b$，定义
\[
 m_i=\inf_{[x_{i-1},x_i]}f,\quad M_i=\sup_{[x_{i-1},x_i]}f,\quad
 L(f,P)=\sum_i m_i\Delta x_i,\quad U(f,P)=\sum_i M_i\Delta x_i,
\]
并称 $\omega_i=M_i-m_i$ 为该区间上的振幅（oscillation）。加细分划使下和不减、上和不增；任意两分划的下和都不超过对方的上和，这可通过共同加细分划证明。因此下积分 $\sup_P L(f,P)$ 不超过上积分 $\inf_P U(f,P)$。

\begin{theorem}
\label{thm:pei-int-darboux}
对有界函数 $f$，下列条件等价：$f$ 可积；上下积分相等；对任意 $\varepsilon>0$ 存在一个分划 $P$ 使
\begin{equation}
 U(f,P)-L(f,P)=\sum_i\omega_i\Delta x_i<\varepsilon;
 \label{eq:pei-int-oscillation-sum}
\end{equation}
以及当 $\lVert P\rVert\to0$ 时，该振幅和趋于零。
\end{theorem}
\begin{proof}
由可积定义，分别选取逼近各小区间上确界、下确界的标记点，可使上和与下和任意接近同一积分值；反过来，积分和总夹在上下和之间。上下积分相等当且仅当存在使其差任意小的分划，这又由共同加细分划及上下积分的确界定义得到。

最后说明“存在一个分划”为何能推出“所有足够细分划”。固定 $P_0$ 使其上下和之差小于 $\varepsilon/2$，设 $P_0$ 有 $r$ 个内分点、$|f|\le M$。任意分划 $Q$ 中，内部不含 $P_0$ 分点的小区间可按 $P_0$ 的区间归组，其振幅和不超过 $P_0$ 的振幅和；其余小区间至多有 $r$ 个，总长不超过 $r\lVert Q\rVert$，贡献至多 $2Mr\lVert Q\rVert$。令网格宽度充分小，就得到 $U(f,Q)-L(f,Q)<\varepsilon$。这也说明两个表述中的量词不同，但在有界条件下确实等价。
\end{proof}

\begin{proposition}
\label{prop:pei-int-large-oscillation}
有界函数 $f$ 可积，当且仅当对每个 $\eta>0$ 和 $\rho>0$，都存在分划，使振幅不小于 $\eta$ 的小区间总长小于 $\rho$。
\end{proposition}
\begin{proof}
必要性由 $\eta\sum_{\omega_i\ge\eta}\Delta x_i\le\sum_i\omega_i\Delta x_i$ 得到。反过来，设 $|f|\le M$，把振幅和分成两类，得 $\sum_i\omega_i\Delta x_i\le\eta(b-a)+2M\rho$；先使 $\eta$ 小，再使 $\rho$ 小，即可应用定理~\ref{thm:pei-int-darboux}。
\end{proof}
连续函数用一致连续性控制每个振幅；单调函数则有 $\sum_i\omega_i\Delta x_i\le\lVert P\rVert|f(b)-f(a)|$。更一般地，若 $f$ 的全变差 $V_a^b(f)=\sup_P\sum_i|f(x_i)-f(x_{i-1})|$ 有限，则每个振幅不超过相应区间的变差，且变差可加，从而振幅和不超过 $\lVert P\rVert V_a^b(f)$。这证明了\emph{有界变差函数（function of bounded variation）}的可积性。

\subsection{把困难点围起来：可积而不连续的函数}
\begin{example}
\label{ex:pei-int-thomae}
沿用例~\ref{ex:pei-thomae} 中的 Thomae 函数 $R$（教材称 Riemann 函数），证明它在 $[0,1]$ 上可积且积分为零。
\end{example}
\begin{solve}
固定 $N$，分母 $q\le N$ 的点只有有限个，用总长小于 $\rho$ 的有限个小区间将它们覆盖，并选分划使覆盖区间的端点成为分点。覆盖区间对上和的贡献至多为 $\rho$，其余部分函数不超过 $1/N$，故 $U(R,P)\le\rho+1/N$；每个非退化区间含无理数，下和恒为零。令 $\rho\downarrow0$、$N\to\infty$，便得积分为零。其连续点的刻画已在例~\ref{ex:pei-thomae} 中证明，这里突出同一函数的可积性。
\end{solve}
同样的方法适用于有界函数的间断点能由总长任意小的有限个区间覆盖的情形：覆盖区间上用有界性控制，余下紧集上用局部连续性和有限覆盖控制振幅。有限个间断点、以及只有有限个聚点的某些间断点列都可这样处理。\emph[caution]{不能把“间断点可数”直接替换成“能被总长任意小的有限个区间覆盖”}；例如上例的有理间断点稠密，证明时须围住函数值较大的有限多个点，而非一次围住全部间断点。

\begin{proposition}
\label{prop:pei-int-composition}
若 $f,g$ 在 $[a,b]$ 上可积，则 $|f|$、$f^2$、$fg$、$\max(f,g)$、$\min(f,g)$ 可积。若 $f([a,b])\subset[A,B]$，而 $\varphi$ 在 $[A,B]$ 上连续，则 $\varphi\circ f$ 可积。
\end{proposition}
\begin{proof}
绝对值满足 $\bigl||u|-|v|\bigr|\le|u-v|$；在有界值域内平方函数满足 Lipschitz 估计。先由振幅判别得到 $|f|$ 和 $f^2$ 的可积性，再用 $fg=((f+g)^2-(f-g)^2)/4$ 与 $\max(f,g)=(f+g+|f-g|)/2$ 等恒等式。一般连续复合的情形，由 $\varphi$ 的一致连续性，对给定 $\varepsilon$ 选择 $\eta$，使 $|u-v|<\eta$ 时 $|\varphi(u)-\varphi(v)|<\varepsilon$；振幅小于 $\eta$ 的区间上复合函数的振幅小，而其余区间总长可以任意小，故由命题~\ref{prop:pei-int-large-oscillation} 得证。
\end{proof}

\begin{notealert}[type=caution,title={不连续的外函数不能直接套用复合结论}]
教材本节单元练习关于取整函数的提示声称取整不增大振幅，这一步一般不成立。取 $f=-R$，其中 $R$ 为例~\ref{ex:pei-int-thomae} 的函数，则 $f$ 可积，而 $\lfloor f(x)\rfloor$ 在有理点等于 $-1$、无理点等于零，每个区间上的振幅均为一，所以不可积。跨过整数阈值时，即使 $f$ 的振幅很小，取整后的振幅仍可为一；使用连续复合定理时必须核对外函数的连续性。
\end{notealert}

\subsection{零积分、连续点与导函数}
\begin{theorem}
\label{thm:pei-int-zero-integral}
可积函数的连续点在 $[a,b]$ 上稠密。对非负可积函数 $f$，$\int_a^bf=0$ 当且仅当 $f$ 在每个连续点均为零。因此连续非负函数的积分为零便处处为零；处处严格为正的可积函数的积分严格为正。
\end{theorem}
\begin{proof}
在任意非退化子区间内，由振幅和可以任意小，能找到一个振幅小于 $1/2$ 的小区间；在其内部取非退化闭区间 $J_1$。重复这一过程，构造 $J_{n+1}\subset\operatorname{int}J_n$，使 $|J_n|\to0$ 且 $f$ 在 $J_n$ 上的振幅小于 $2^{-n}$。区间套的公共点 $c$ 对每个 $J_n$ 都是内点，故 $f$ 在 $c$ 连续。初始子区间任意，因而连续点稠密。

若连续点 $c$ 满足 $f(c)>0$，则某个邻域内 $f\ge f(c)/2$，积分必为正。反过来，若所有连续点处 $f=0$，每个非退化小区间都有一个这样的点，所以每个 Darboux 下和为零，由可积性得积分为零。最后两条结论分别利用处处连续和至少存在一个连续点即得。
\end{proof}
若连续函数 $f\ge g$ 且二者积分相等，则对连续非负函数 $f-g$ 应用上述结论，得到 $f\equiv g$。积分估计中一串不等式首尾相等时，这一事实可逐步迫使被积函数相等，进而与可微性等额外条件矛盾。

\begin{proposition}
设 $F\in C^1[a,b]$，则 $F'$ 可积且 $\int_a^bF'=F(b)-F(a)$。更一般地，设 $F$ 在 $[a,b]$ 连续、在 $(a,b)$ 可微，并对导数作端点赋值，则 $F'$ 可积当且仅当存在可积函数 $g$ 使 $F(x)=F(a)+\int_a^xg(t)\,dt$。
\end{proposition}
\begin{proof}
第一条既可由连续可积性和微积分基本定理得到，也可逐段用微分中值定理，将 $F(b)-F(a)$ 写成 $\sum_i F'(\eta_i)\Delta x_i$，再用 $F'$ 的一致连续性与任意标记的和比较。

对第二条，必要性取 $g=F'$。充分性中，固定一个小区间 $J$，记 $m=\inf_Jg$、$M=\sup_Jg$；由积分平均值的上下界，对其内任意 $u<v$ 有 $m\le(F(v)-F(u))/(v-u)\le M$，令端点向任意内点逼近，得到 $m\le F'\le M$。故除有限分点的赋值外，$F'$ 的振幅受 $g$ 的相应振幅控制，由 Darboux 判别推出可积性；有限个点的赋值不影响可积性与积分。不能直接将积分表达式求导后声称处处 $F'=g$，因为 $g$ 未必处处连续。
\end{proof}

\subsection{积分定义可以怎样简化}
定义中的两种任意性不能同时随意删掉。对有界 $f$，以下两种检验仍足够：其一，只取等分分划，但要求\emph{所有标记点的积分和一致趋于同一极限}；其二，保留任意分划，每个小区间只选左、右端点之一，但要求对所有这些选择统一收敛。第一种检验可由等分分划的上下和证明；第二种检验中，把任意标记点所在的小区间在标记点处一分为二，再在两个新小区间都选该标记点，所得端点和恰好等于原标记和，且网格不增大，故回到原定义。若既只取等分分划又只取右端点，条件就不够：$[0,1]$ 上的有理数示性函数的所有等分右端点和恒为一，但函数不可积。等价地，可积性还可写成积分和的 Cauchy 条件，即任意两个充分细分划、任意两组标记点所给的积分和之差均任意小；其充分性利用实数完备性先取得一个极限，再与任意积分和比较。

\section{有关积分的几类典型问题}
\label{sec:pei-int-integral-problems}
% 教材对应：§4.3，PDF 第 33--82 页（原书第 282--331 页）。

\subsection{从已知性质建立积分上下界}
积分估计首先问“已知信息能给出什么逐点界”，再问是否需要换元、分部积分或分区间。若 $f$ 单调递增，等分步长为 $h=(b-a)/n$，则左、右端点和夹住积分，且二者之差恰为 $h(f(b)-f(a))$；因此指定精度时，可先选 $n$ 使这个差小于容许误差，再计算有限和。一般连续函数也可用 Darboux 上下和逼近，但仅知连续性时不能直接套用这一单调误差公式。

\begin{proposition}
\label{prop:pei-int-midpoint-bound}
若 $f\in C^2[a,b]$，$c=(a+b)/2$，$|f''|\le M$，则
\[
 \left|\int_a^bf(x)\,dx-(b-a)f(c)\right|\le\frac{M(b-a)^3}{24}.
\]
还存在 $\xi\in(a,b)$，使差恰等于 $f''(\xi)(b-a)^3/24$。
\end{proposition}
\begin{proof}
在 $c$ 处应用带积分余项的 Taylor 公式。线性项在对称区间上积分为零，二阶余项的绝对值不超过 $M(x-c)^2/2$，故误差不超过 $M(b-a)^3/24$。误差还可写成 $f''$ 乘以一个非负核的积分：左半区间的核为 $(t-a)^2/2$，右半区间的核为 $(b-t)^2/2$。核积分为 $(b-a)^3/24$，且在内部为正；由加权积分中值定理得到 $\xi$。
\end{proof}

\begin{example}
证明 $\displaystyle\int_0^{\sqrt{2\pi}}\sin(x^2)\,dx>0$。
\end{example}
\begin{solve}
令 $t=x^2$，并把 $[0,2\pi]$ 分成两个半周期，在第二段再令 $t=u+\pi$，得到
\[
 \int_0^{\sqrt{2\pi}}\sin(x^2)\,dx
 =\frac12\int_0^\pi\sin t\left(\frac1{\sqrt t}-\frac1{\sqrt{t+\pi}}\right)\,dt>0.
\]
零点附近的反常积分因 $t^{-1/2}\sin t=O(t^{1/2})$ 而收敛；括号在内部严格为正。换元把非线性相位变成等长半周期，配对后利用权函数递减性，便可判断正负而不求原函数。
\end{solve}
反复分部积分的作用是把导数从难估计的函数转到可控的权函数上。例如 $f,g\in C^{2m}[a,b]$，且 $f^{(j)}$、$g^{(j)}$ 在两端均为零（$0\le j<m$），则所有边界项消失，得到 $\int_a^bf^{(2m)}g=\int_a^bfg^{(2m)}$。选择 $g(x)=(x-a)^m(b-x)^m\ge0$，其最高阶导数为常数，便能把 $\int f$ 与最高阶导数的界联系起来。使用这一方法时应逐个检查边界项，不能仅因函数值为零就删去含高阶导数的边界项。

\subsection{导数约束、最佳常数与等号检查}
\begin{proposition}
\label{prop:pei-int-endpoint-slope}
若 $f\in C^1[a,b]$、$f(a)=f(b)=0$，则
\begin{equation}
 \int_a^b|f(x)|\,dx\le\frac{(b-a)^2}{4}\max_{[a,b]}|f'|.
 \label{eq:pei-int-slope-area}
\end{equation}
若 $f$ 非零，则在上述 $C^1$ 条件下不等号严格，但常数 $1/4$ 仍不能减小。
\end{proposition}
\begin{proof}
记 $D=\max|f'|$，由两端点分别应用中值定理，$|f(x)|\le D\min(x-a,b-x)$。积分后即得式~\eqref{eq:pei-int-slope-area}。若 $D>0$ 且取等号，由连续非负函数零积分的性质，必有 $|f(x)|=D\min(x-a,b-x)$；这要求中点处出现尖角，与 $C^1$ 矛盾。另一方面，把折线 $\min(x-a,b-x)$ 在中点的任意小邻域内光滑化，可使导数上界保持在一以内、面积任意接近 $(b-a)^2/4$，因此常数不能减小。
\end{proof}
\emph[caution]{最佳常数不意味着存在取等号的光滑函数}。例如若要求 $f\in C^1[0,2]$、$f(0)=f(2)=1$、$|f'|\le1$，则必有 $f(x)\ge|x-1|$，从而 $\int_0^2f\ge1$；若还要求积分恰为一，则连续性强制 $f=|x-1|$，与可微性矛盾。积分估计不仅给数值范围，也能排除函数的存在性。

\subsection{积分不等式中的单调性、对称性与求导}
若 $f$ 在 $[0,1]$ 上单调不增，则对 $0<c<1$，前段平均值不小于后段平均值，即
\[
 \frac1c\int_0^cf\ge\frac1{1-c}\int_c^1f,
 \qquad \int_0^cf\ge c\int_0^1f.
\]
可将 $f(x)\ge f(y)$ 在 $0\le x\le c\le y\le1$ 上先后积分得到，故不要求 $f$ 非负。若 $f$ 在 $[a,b]$ 单调递增，$c=(a+b)/2$，则
\[
 \int_a^b(x-c)f(x)\,dx
 =\int_0^{(b-a)/2}t\bigl(f(c+t)-f(c-t)\bigr)\,dt\ge0.
\]
对称点配对常能把两项的差变成逐点非负表达式。类似地，若 $f$ 在 $[-a,a]$ 上连续且为偶函数，则由 $x\mapsto-x$ 和 $1/(1+e^x)+1/(1+e^{-x})=1$ 得到 $\int_{-a}^a f(x)/(1+e^x)\,dx=\int_0^af(x)\,dx$。

\begin{example}
设 $f\in C^1[0,1]$、$f(0)=0$、$0<f'(x)<1$（$0<x<1$）。证明 $\displaystyle\left(\int_0^1f\right)^2>\int_0^1f^3$。
\end{example}
\begin{solve}
令 $A(x)=\int_0^xf(t)\,dt$，$H(x)=2A(x)-f(x)^2$。有 $H(0)=0$、$H'=2f(1-f')>0$，故 $H(x)>0$。再令 $F(x)=A(x)^2-\int_0^xf(t)^3\,dt$，则 $F(0)=0$、$F'=fH>0$，从而 $F(1)>0$。把常上限临时改成变量上限，再构造两端之差，是用微分学证明积分不等式的基本办法。
\end{solve}
从已知逐点不等式出发反复积分，也能生成更高阶的估计。例如由 $\cos x\le1$（$x\ge0$），依次积分得到 $\sin x\le x$、$1-\cos x\le x^2/2$，继而得到 $x-x^3/6\le\sin x\le x-x^3/6+x^5/120$。严格性须检查初始不等式是否在具有正长度的区间上严格成立。

\subsection{正交条件、零点与线性相关}
\begin{proposition}
\label{prop:pei-int-moment-zeros}
设非零连续函数 $f$ 在 $[a,b]$ 上满足 $\int_a^bx^kf(x)\,dx=0$（$k=0,\ldots,m-1$）。则 $f$ 在 $(a,b)$ 内至少有 $m$ 个不同零点。
\end{proposition}
\begin{proof}
若零点已经有无穷多个，则结论成立；否则假设内部零点少于 $m$，从中挑出真正改变符号的点 $c_1,\ldots,c_r$，$r<m$。多项式 $p(x)=\prod_{j=1}^r(x-c_j)$（$r=0$ 时取一）与 $f$ 在每个变号点同时变号，故 $pf$ 在整个区间上保持同号，且在某个非退化小区间上严格非零。因此 $\int pf\ne0$；但 $\deg p<m$，由矩条件又有 $\int pf=0$，矛盾。
\end{proof}
“选一个在零点处同步变号的乘子”是把多个积分等式转为零点数量的通用方法。例如 $\int_0^1f=\int_0^1xf=0$ 保证至少两个内部零点；仅分别对两个积分套一次中值定理，不能保证所得零点不同。

\begin{proposition}
\label{prop:pei-int-gram}
连续函数 $f_1,\ldots,f_m$ 在 $[a,b]$ 线性相关，当且仅当 Gram 矩阵 $G=(\int_a^bf_i f_j)_{i,j=1}^m$ 的行列式为零。
\end{proposition}
\begin{proof}
对 $c\in\mathbb R^m$，$c^{\mathsf T}Gc=\int_a^b(\sum_i c_if_i)^2\ge0$，所以 $G$ 半正定。若函数线性相关，相应非零系数向量属于 $G$ 的核。反过来，$\det G=0$ 时取非零核向量 $c$，其平方积分为零；被积函数连续非负，故 $\sum_i c_if_i\equiv0$。仅有可积性时，平方积分为零不能推出每一点上都为零，连续性在这里确有作用。
\end{proof}
积分还能把覆盖重数转为长度估计：若 $[0,1]$ 被 $n$ 个区间覆盖至少 $q$ 次，区间示性函数之和处处不小于 $q$，积分后得到区间总长不小于 $q$，故至少一个区间的长度不小于 $q/n$。这说明合适的示性函数可以把几何计数问题变成积分的可加性与保序性。

\subsection{凸性与积分平均}
\begin{theorem}
\label{thm:pei-int-hadamard-jensen}
若 $f$ 在 $[a,b]$ 上连续且为凸函数（convex function），则有 Hermite--Hadamard 不等式
\[
 f\left(\frac{a+b}{2}\right)\le\frac1{b-a}\int_a^bf(x)\,dx\le\frac{f(a)+f(b)}2.
\]
若 $w\ge0$ 可积、$W=\int_a^bw>0$，$u$ 可积且取值于 $[A,B]$，$\varphi$ 在 $[A,B]$ 连续且为凸函数，则加权 Jensen 不等式为
\begin{equation}
 \varphi\left(\frac1W\int_a^bwu\right)\le\frac1W\int_a^bw\varphi(u).
 \label{eq:pei-int-jensen}
\end{equation}
\end{theorem}
\begin{proof}
对称点 $x,a+b-x$ 的凸性不等式给出左侧，连接两端点的弦位于函数图像上方，积分给出右侧。Jensen 不等式令 $c=W^{-1}\int wu$；若 $A<c<B$，在 $c$ 处取一条支撑直线 $\ell(t)=\varphi(c)+s(t-c)\le\varphi(t)$，乘以 $w$ 后积分，线性项消失，便得结论。若 $c=A$，则 $\int w(u-A)=0$。对任意 $\varepsilon>0$，连续性给出 $\delta>0$，而有界性给出常数 $K$，使处处 $|\varphi(u)-\varphi(A)|\le\varepsilon+(K/\delta)(u-A)$；乘以 $w$ 积分后误差不超过 $\varepsilon W$，令 $\varepsilon\downarrow0$ 即得等号。$c=B$ 时同理，$A=B$ 时直接成立。若 $u,w$ 均连续，且 $w>0$、$\varphi$ 严格凸，等号当且仅当 $u$ 为常数。
\end{proof}
若 $f$ 在 $[0,\infty)$ 连续且凸，$A(x)=x^{-1}\int_0^xf(t)\,dt$（$x>0$），则换元给出 $A(x)=\int_0^1f(tx)\,dt$，直接对被积函数应用凸性就得到 $A$ 也凸；不必先求二阶导数。若 $g$ 单调递增，则 $F(x)=\int_c^xg(t)\,dt$ 为凸函数，因为相邻区间上的平均值随区间向右移动而不减，即割线斜率递增。这两个结论展示了“积分保持凸性”和“单调函数积分生成凸函数”的不同方向。

\subsection{变上限积分、函数方程与积分不等式}
若 $f$ 局部 Riemann 可积，$F(x)=\int_c^xf(t)\,dt$ 总是连续；在 $f$ 连续的点才可直接断言 $F'(x)=f(x)$。若 $f$ 连续、$u,v$ 可微，则
\[
 \frac{d}{dx}\int_{u(x)}^{v(x)}f(t)\,dt=f(v(x))v'(x)-f(u(x))u'(x).
\]
该公式先用一个固定基点的原函数拆成两项，再应用链式法则，负向积分亦自动包含其中。比如对局部可积的加性函数 $f(x+y)=f(x)+f(y)$，在 $y\in[0,1]$ 上积分得
\[
 f(x)=\int_x^{x+1}f(t)\,dt-\int_0^1f(t)\,dt,
\]
右端连续，从而 $f$ 连续。再由加性在有理点上确定 $f(r)=rf(1)$，用连续性延拓到实数，得到 $f(x)=xf(1)$。这里积分首先提供正则性，然后才使用连续加性函数的结论。

\begin{proposition}
\label{prop:pei-int-gronwall}
设 $u,v\ge0$ 在 $[a,b]$ 上连续，$c\ge0$，且 $u(x)\le c+\int_a^xv(t)u(t)\,dt$，则
\[
 u(x)\le c\exp\left(\int_a^xv(t)\,dt\right).
\]
\end{proposition}
\begin{proof}
记 $H(x)=c+\int_a^xvu$，则 $H'\le vH$，故 $\bigl(e^{-\int_a^xv}H(x)\bigr)'\le0$，从而 $H(x)\le ce^{\int_a^xv}$，再用 $u\le H$。当 $c=0$ 时相同论证给出 $u\equiv0$，无需除以 $u$。
\end{proof}
这就是 Gronwall 不等式的一个基本形式。教材中含未知函数积分的不等式，常可先令该积分为辅助函数，再乘积分因子；对正函数的微分不等式，也可积分其对数导数，例如 $f'/f\le-\alpha/x$、$f>0$ 推出 $f(x)\le f(a)(a/x)^\alpha$。非负函数可能有零点时，不应未经处理就除以 $f$。

\subsection{第一积分中值定理及其内点形式}
\begin{theorem}
\label{thm:pei-int-first-mean}
设 $f$ 在 $[a,b]$ 连续，$g\ge0$ 可积，则存在 $\xi\in(a,b)$，使
\[
 \int_a^bf(x)g(x)\,dx=f(\xi)\int_a^bg(x)\,dx.
\]
$g\le0$ 时同样成立；特别地 $g\equiv1$ 时，积分平均值能在区间内部取得。
\end{theorem}
\begin{proof}
若 $\int g=0$，由 $|\int fg|\le\max|f|\int g=0$，任意内点均可。否则记加权平均值为 $r$，上下界保证 $\min f\le r\le\max f$。若内部没有 $f=r$ 的点，连续性迫使 $f-r$ 在整个 $(a,b)$ 严格同号。又因 $g$ 的积分为正，可以截去足够短的两端后仍保留正积分，而 $|f-r|$ 在这个内闭区间有正下界，故 $\int(f-r)g\ne0$，矛盾。
\end{proof}
定理要求权函数不变号；若 $g$ 变号，积分平均值可能落在 $f$ 的值域之外。对于仅有界可积而无介值性的 $f$，通常只能得到 $\int fg=\theta\int g$、$\inf f\le\theta\le\sup f$，不能把 $\theta$ 写成某个函数值。强化为内点在多次使用 Rolle 定理时有实际意义：分别在互不相交的开子区间内取得中值点，才能确保它们不同，随后逐次计数导数的零点。

\subsection{第二积分中值定理与振荡估计}
\begin{theorem}
\label{thm:pei-int-second-mean}
设 $g$ 在 $[a,b]$ 可积。若 $f$ 非负单调不增，则存在 $\xi\in[a,b]$ 使 $\int_a^bfg=f(a)\int_a^\xi g$；若 $f$ 非负单调不减，则存在 $\xi$ 使 $\int_a^bfg=f(b)\int_\xi^bg$。一般地，只要 $f$ 单调，就有某个 $\xi\in[a,b]$ 满足 Bonnet 公式
\begin{equation}
 \int_a^bfg=f(a)\int_a^\xi g+f(b)\int_\xi^bg.
 \label{eq:pei-int-bonnet}
\end{equation}
\end{theorem}
\begin{proof}
先设 $f\ge0$ 单调不增，令 $G(x)=\int_a^xg$。在分划上用右端点阶梯函数逼近 $f$，其与 $f$ 的积分误差不超过 $\sup|g|\lVert P\rVert(f(a)-f(b))$。相应近似积分经 Abel 变换为
\[
 f(x_m)G(b)+\sum_{i=1}^{m-1}\bigl(f(x_i)-f(x_{i+1})\bigr)G(x_i).
\]
所有系数非负，总和为 $f(x_1)\le f(a)$；补上 $G(a)=0$ 的系数后，总和恰为 $f(a)$。故其值夹在 $f(a)\min G$ 与 $f(a)\max G$ 之间。取极限，利用 $G$ 连续及介值性即得第一种形式；递增的情形反射区间即可。一般递减的 $f$ 对 $f-f(b)\ge0$ 应用第一种形式，移项得到式~\eqref{eq:pei-int-bonnet}；递增的情形同理。
\end{proof}
例如 $f\ge0$ 递减时，$|\int_a^bf(x)\sin x\,dx|\le2f(a)$。若相位 $\varphi\in C^1[a,b]$，$\varphi'\ge m>0$ 且 $\varphi'$ 单调，令 $t=\varphi(x)$ 后，振幅 $1/\varphi'(\varphi^{-1}(t))$ 非负单调，故
\[
 \left|\int_a^be^{\mathrm i\varphi(x)}\,dx\right|\le\frac2m.
\]
这里复指数只是同时记述正弦和余弦；同一 Abel 变换证明对复数积分给出上述模长界，因为 $|\int_\alpha^\beta e^{\mathrm it}\,dt|\le2$。若只讨论实数，也可分别得到正弦和余弦积分的同样界。第一中值定理控制不变号的权，第二中值定理利用单调振幅与原函数的有界性，两者的假设和用途应分别核对。

\section{几个著名的不等式}
\label{sec:pei-int-classical-inequalities}
% 教材对应：§4.4，PDF 第 82--102 页（原书第 331--351 页）。

\subsection{Cauchy--Schwarz 不等式与等号条件}
\begin{theorem}
\label{thm:pei-int-schwarz}
对实数列 $(a_i)$、$(b_i)$，有 $|\sum_i a_ib_i|^2\le(\sum_i a_i^2)(\sum_i b_i^2)$，等号当且仅当两个向量线性相关。若 $f,g$ 在 $[a,b]$ Riemann 可积，则积分形式为
\begin{equation}
 \left|\int_a^bfg\right|^2\le\left(\int_a^bf^2\right)\left(\int_a^bg^2\right).
 \label{eq:pei-int-schwarz}
\end{equation}
一般可积函数取等号的充要条件是存在不全为零的 $\alpha,\beta$ 使 $\int_a^b(\alpha f+\beta g)^2=0$；若二者连续，则等价于 $\alpha f+\beta g\equiv0$。
\end{theorem}
\begin{proof}
二次式 $\int(f+tg)^2=A+2Bt+Ct^2$ 对每个实数 $t$ 非负。当 $C>0$ 时，其最小值 $A-B^2/C\ge0$，即得不等式，且等号恰对应某个非零线性组合的平方积分为零。若 $C=0$，对所有 $t$ 非负迫使 $B=0$，不等式也成立。离散形式把积分换成有限求和，证明相同；连续情形的等号条件由定理~\ref{thm:pei-int-zero-integral} 得到。
\end{proof}
命题~\ref{prop:pei-int-gram} 给出了多函数推广：Gram 矩阵半正定，其行列式非负，连续函数情形下行列式为零恰对应线性相关。应用 Schwarz 不等式时，选择分拆常比公式本身更关键：把 $f=\sqrt w\,(f/\sqrt w)$ 可得 $|\int f|^2\le(\int w)(\int f^2/w)$，其中 $w$ 应严格为正，所有积分须存在。

\begin{example}
设 $g\in C^1[a,b]$、$g(a)=0$，证明
\[
 \int_a^b|gg'|\le\frac{b-a}{2}\int_a^b(g')^2,
\]
并确定等号条件。若还要求 $g(b)=0$，如何改进常数？
\end{example}
\begin{solve}
令 $H(x)=\int_a^x|g'(t)|\,dt$，则 $|g|\le H$，所以
\[
 \int_a^b|gg'|\le\int_a^bHH'=\frac12\left(\int_a^b|g'|\right)^2
 \le\frac{b-a}{2}\int_a^b(g')^2.
\]
取等号要求 $|g'|$ 为常数；$g'$ 连续，非零时不能改变符号，故 $g(x)=c(x-a)$，反过来这些函数均取等号。若两端为零，在中点处分开，分别从左右端应用刚证的单端点结论，得到 $\int|gg'|\le(b-a)\int(g')^2/4$。非零 $C^1$ 函数不能取等号，但将两段直线构成的帐篷函数在顶点附近光滑化，可说明常数 $(b-a)/4$ 不能减小。
\end{solve}
另一个常用推论是 $g(a)=0$ 时 $\int_a^bg^2\le(b-a)^2\int_a^b(g')^2/2$：由 $g(x)=\int_a^xg'$ 得到 $g(x)^2\le(x-a)\int_a^b(g')^2$，再积分即可。它提供一个方便的估计，此处并不声称常数最佳；重复应用可把低阶导数的积分控制在高阶导数的积分之下，但每一步都必须有相应的端点消失条件。

\subsection{加权平均与幂平均的积分形式}
设 $w>0$、$f>0$ 在 $[a,b]$ 上连续，定义归一化平均 $\mathcal A_w(u)=(\int wu)/(\int w)$，以及\emph{幂平均（power mean）}与\emph{几何平均（geometric mean）}
\[
 M_r(f;w)=\bigl(\mathcal A_w(f^r)\bigr)^{1/r}\quad(r\ne0),
 \qquad M_0(f;w)=\exp\bigl(\mathcal A_w(\ln f)\bigr).
\]
归一化保证 $M_r(c;w)=c$。在紧区间上 $f$ 有正下界，因此负次幂与对数均可积。

\begin{theorem}
\label{thm:pei-int-power-means}
对任意实数 $r<s$，$M_r(f;w)\le M_s(f;w)$；若 $f$ 非常数，则不等号严格。并且 $\lim_{r\to0}M_r(f;w)=M_0(f;w)$，$\lim_{r\to\infty}M_r(f;w)=\max f$。
\end{theorem}
\begin{proof}
当 $0<r<s$ 时，对 $f^r$ 和凸函数 $t^{s/r}$ 应用式~\eqref{eq:pei-int-jensen}，得到 $\mathcal A_w(f^s)\ge\mathcal A_w(f^r)^{s/r}$。当 $r<s<0$ 时，对 $1/f$ 应用刚证的正指数结论并取倒数。当指数跨过零时，先由凸函数 $e^{rt}$ 的 Jensen 不等式得到 $M_r\ge M_0$（$r>0$），并对 $1/f$ 得到 $M_r\le M_0$（$r<0$）。严格性来自相应函数的严格凸性与 $w>0$。

由于 $\ln f$ 有界，$f^r=1+r\ln f+O(r^2)$ 的余项一致，故 $r^{-1}\ln\mathcal A_w(f^r)\to\mathcal A_w(\ln f)$，得到零点处的极限。无穷大处的极限由定理~\ref{thm:pei-int-power-integral} 得到，归一化常数的 $r$ 次方根趋于一。
\end{proof}
特别地，$M_{-1}\le M_0\le M_1\le M_2$ 分别对应调和、几何、算术与平方平均的次序。离散加权平均把 $\mathcal A_w(u)$ 换成 $\sum_i p_i u_i$（$p_i>0$、$\sum_i p_i=1$），同一 Jensen 证明仍适用；非归一化权重须先除以总权重。若只讨论正指数，允许 $f$ 为零，但对数、负指数及等号条件需另外处理。

\subsection{Hölder 不等式及反向形式}
\begin{theorem}
\label{thm:pei-int-holder}
设 $p,q>1$、$1/p+1/q=1$，$f,g$ 在 $[a,b]$ 可积，则
\begin{equation}
 \int_a^b|fg|\le\left(\int_a^b|f|^p\right)^{1/p}
                     \left(\int_a^b|g|^q\right)^{1/q}.
 \label{eq:pei-int-holder}
\end{equation}
若 $f,g$ 连续且两个幂积分均非零，等号当且仅当 $|f|^p=C|g|^q$，其中 $C>0$。若 $0<p<1$、$q=p/(p-1)<0$，且 $f\ge0$、$g>0$ 连续，则反向形式为
\[
 \int_a^bfg\ge\left(\int_a^bf^p\right)^{1/p}\left(\int_a^bg^q\right)^{1/q}.
\]
在 $f$ 不恒为零时，等号仍以 $f^p$ 与 $g^q$ 成比例刻画。
\end{theorem}
\begin{proof}
先证普通形式。令 $A=(\int|f|^p)^{1/p}$、$B=(\int|g|^q)^{1/q}$；$AB=0$ 时由下段的直接估计处理。否则对 $u=|f|/A$、$v=|g|/B$ 使用数值 Young 不等式 $uv\le u^p/p+v^q/q$，积分得 $\int uv\le1$。连续情形等号要求逐点 $u^p=v^q$。反向形式则在普通 Hölder 不等式中使用共轭指数 $1/p$、$1/(1-p)$：
\[
 \int f^p=\int(fg)^p(g^q)^{1-p}
 \le(\int fg)^p(\int g^q)^{1-p},
\]
移项即得结论。由于 $q<0$，$g$ 的正下界是保证各项有意义的一个充分条件。
\end{proof}
证明中若 $AB=0$，也可完全用 Riemann 积分估计处理：例如 $\int|f|^p=0$，对任意 $\delta>0$，在 $|f|\le\delta$ 时有 $|f|\le\delta$，其余情形有 $|f|\le\delta^{1-p}|f|^p$，故处处 $|f|\le\delta+\delta^{1-p}|f|^p$。积分后令 $\delta\downarrow0$，得到 $\int|f|=0$，再用 $g$ 有界得到 $\int|fg|=0$。把积分换成求和可得离散形式；在积分中改用归一化权重，则得到加权形式。$p=q=2$ 时回到 Schwarz 不等式。

\subsection{Minkowski 不等式与参数范围}
\begin{theorem}
\label{thm:pei-int-minkowski}
对 $p\ge1$，记 $\lVert f\rVert_p=(\int_a^b|f|^p)^{1/p}$，则
\begin{equation}
 \lVert f+g\rVert_p\le\lVert f\rVert_p+\lVert g\rVert_p.
 \label{eq:pei-int-minkowski}
\end{equation}
若 $f,g\ge0$ 连续、$p>1$ 且均不恒为零，等号当且仅当二者成比例。若 $0<r<1$、$f,g\ge0$ 连续，则
\[
 \left(\int_a^b(f+g)^r\right)^{1/r}
 \ge\left(\int_a^bf^r\right)^{1/r}+\left(\int_a^bg^r\right)^{1/r}.
\]
$r<0$ 时，在 $f,g>0$ 连续的条件下也有后一反向不等式。
\end{theorem}
\begin{proof}
$p=1$ 由逐点三角不等式得到。$p>1$ 时，令 $q=p/(p-1)$，利用 $|f+g|\le|f|+|g|$ 与 Hölder 不等式，
\[
 \int|f+g|^p\le\int(|f|+|g|)|f+g|^{p-1}
 \le(\lVert f\rVert_p+\lVert g\rVert_p)\lVert f+g\rVert_p^{p-1}.
\]
总范数为零时直接成立，否则约去相应幂次。非负连续情形由 Hölder 的等号条件得到比例关系。

对反向形式，设 $A=(\int f^r)^{1/r}$、$B=(\int g^r)^{1/r}$ 均为正，令 $\lambda=A/(A+B)$、$u=f/A$、$v=g/B$。当 $0<r<1$ 时，$t^r$ 凹，故 $\int(\lambda u+(1-\lambda)v)^r\ge1$；当 $r<0$ 时，$t^r$ 凸，积分不等号相反，但取负指数 $1/r$ 又反向。两种情况都得到 $\lVert f+g\rVert_r\ge A+B$。$0<r<1$ 时某函数恒为零的情形直接成立。
\end{proof}
\emph[caution]{在 $p=1$ 时，非负函数总有 $\lVert f+g\rVert_1=\lVert f\rVert_1+\lVert g\rVert_1$，无需成比例}；只有 $p>1$ 才能这样使用严格凸性给出的等号条件。对一般 Riemann 可积函数，$\lVert f\rVert_p=0$ 不保证逐点为零，因此 $p\ge1$ 时该表达式在原函数空间上是半范数（seminorm）；限制为连续函数，或将满足 $\int|f-g|^p=0$ 的函数视为同一个对象后，才得到范数。$0<r<1$ 时的表达式不满足通常的三角不等式。多函数形式可由二函数形式归纳得到，离散形式同样由求和替代积分。

\subsection{Young 不等式与反函数的面积恒等式}
\begin{theorem}
\label{thm:pei-int-young}
设 $\phi:[0,\infty)\to[0,\infty)$ 连续、严格递增、满射，且 $\phi(0)=0$。则对 $a,b\ge0$，
\begin{equation}
 ab\le\int_0^a\phi(x)\,dx+\int_0^b\phi^{-1}(y)\,dy,
 \label{eq:pei-int-young-inverse}
\end{equation}
等号当且仅当 $b=\phi(a)$。取 $\phi(x)=x^{p-1}$（$p>1$）即得 $ab\le a^p/p+b^q/q$，其中 $q=p/(p-1)$，等号当且仅当 $a^p=b^q$。
\end{theorem}
\begin{proof}
先用分划 $0=x_0<\cdots<x_n=a$ 与像分划 $\phi(x_i)$ 的积分和，或对矩形作曲边区域分割，得到
\[
 \int_0^a\phi(x)\,dx+\int_0^{\phi(a)}\phi^{-1}(y)\,dy=a\phi(a).
\]
固定 $a$ 后，式~\eqref{eq:pei-int-young-inverse} 右端减去左端等于 $\int_{\phi(a)}^b(\phi^{-1}(y)-a)\,dy$。若 $b>\phi(a)$，被积函数在内部为正；若 $b<\phi(a)$，积分方向与被积函数的符号同时反向；故均为正，只有 $b=\phi(a)$ 时为零。这一证明不要求 $\phi$ 可微。
\end{proof}

\subsection{Chebyshev 不等式：同序与反序}
\begin{proposition}
\label{prop:pei-int-chebyshev}
若可积函数 $f,g$ 同序，即 $(f(x)-f(y))(g(x)-g(y))\ge0$ 对所有 $x,y\in[a,b]$ 成立，$w\ge0$ 可积且 $W=\int w>0$，则
\[
 \left(\int wf\right)\left(\int wg\right)\le W\int wfg.
\]
反序时不等号相反。特别地，同为递增或同为递减的函数同序，取 $w\equiv1$ 得无权形式。
\end{proposition}
\begin{proof}
展开并分别积分可得
\[
 W\int wfg-\int wf\int wg
 =\frac12\int_a^b\int_a^b w(x)w(y)(f(x)-f(y))(g(x)-g(y))\,dx\,dy.
\]
由同序条件，右端非负。离散形式的差同样等于 $\tfrac12\sum_{i,j}p_ip_j(a_i-a_j)(b_i-b_j)$，因此加权推广与无权形式具有同一证明。
\end{proof}

\section{反常积分}
\label{sec:pei-int-improper-integrals}
% 教材对应：§4.5，PDF 第 102--136 页（原书第 351--385 页）。

\subsection{先识别反常点，再作计算}
\begin{definition}
若 $f$ 在每个 $[a,R]$ 上 Riemann 可积，且有限极限 $\lim_{R\to\infty}\int_a^Rf$ 存在，称\emph{反常积分（improper integral）}$\int_a^\infty f$ 收敛。若 $f$ 在每个 $[c,b]$ 上可积，则 $\int_a^bf=\lim_{c\downarrow a}\int_c^bf$ 定义端点无界情形。存在多个反常点时，必须分开每个反常点并分别取极限；各部分都收敛，才称整个反常积分收敛。若 $\int|f|<\infty$，称绝对收敛；收敛而不绝对收敛，称条件收敛。
\end{definition}
\emph[caution]{对称截断的 Cauchy 主值（Cauchy principal value）不等于通常的反常积分}。例如 $\int_{-1}^1 dx/x$ 在零点两侧均发散，虽然对称删去 $(-\varepsilon,\varepsilon)$ 后的积分恒为零，其主值为零。变量替换、分部积分和原函数计算应先在没有反常点的截断区间上进行，再分别检查边界极限；不能依靠两个发散边界项相减得到积分值。

\begin{example}
计算 $\displaystyle I=\int_0^\infty\frac{dx}{1+x^4}$。
\end{example}
\begin{solve}
被积函数在零点附近有界，在无穷远为 $O(x^{-4})$，故绝对收敛。换元 $x=1/t$ 得 $I=\int_0^\infty x^2/(1+x^4)\,dx$，相加后再令 $u=x-1/x$，得到
\[
 2I=\int_0^\infty\frac{1+x^2}{1+x^4}\,dx
 =\int_{-\infty}^\infty\frac{du}{u^2+2}=\frac\pi{\sqrt2},
 \qquad I=\frac\pi{2\sqrt2}.
\]
倒数变换先制造互补的分子，再用双射换元消去四次式；各步都可由有限截断后取极限验证。
\end{solve}

\begin{proposition}
设 $a,b>0$，$F:[2\sqrt{ab},\infty)\to\mathbb R$ 连续。若 $F\ge0$（允许积分值为 $+\infty$），或所讨论的积分绝对收敛，则
\[
 \int_0^\infty F\left(ax+\frac bx\right)\,dx
 =\frac1a\int_0^\infty F\left(\sqrt{t^2+4ab}\right)\,dt.
\]
因此，以已知 Gaussian 积分 $\int_0^\infty e^{-t^2}\,dt=\sqrt\pi/2$ 为前提，有
\[
 \int_0^\infty\exp\left(-a^2x^2-\frac{b^2}{x^2}\right)\,dx
 =\frac{\sqrt\pi}{2a}e^{-2ab}.
\]
\end{proposition}
\begin{proof}
令 $t=ax-b/x$，它把 $(0,\infty)$ 严格递增地映到 $\mathbb R$，且
\[
 ax+\frac bx=\sqrt{t^2+4ab},\qquad
 \frac{dx}{dt}=\frac1{2a}\left(1+\frac t{\sqrt{t^2+4ab}}\right).
\]
将 $t>0$ 与 $t<0$ 两部分配对，奇函数项相消，便得第一式；对非负 $F$ 可先配对有限区间再取单调极限，绝对收敛时也允许分项。取 $F(u)=e^{-u^2}$，再利用 $(ax+b/x)^2=a^2x^2+b^2/x^2+2ab$，得到第二式。
\end{proof}

\subsection{递推、对称与自我消去}
\begin{proposition}
对 $\alpha>-1$、整数 $m\ge0$，有
\[
 \int_0^1x^\alpha(\ln x)^m\,dx=\frac{(-1)^m m!}{(\alpha+1)^{m+1}}.
\]
对非负整数 $m,n$，有 $\int_0^1x^m(1-x)^n\,dx=m!n!/(m+n+1)!$。
\end{proposition}
\begin{proof}
第一式在 $[\varepsilon,1]$ 上分部积分，边界项因 $\varepsilon^{\alpha+1}|\ln\varepsilon|^m\to0$ 而消失，得到 $I_m=-mI_{m-1}/(\alpha+1)$，从 $I_0=1/(\alpha+1)$ 递推。第二式同样分部积分，得到 $B_{m,n}=nB_{m+1,n-1}/(m+1)$，递推到 $B_{m+n,0}=1/(m+n+1)$ 即得。
\end{proof}
Wallis 积分的初值、降阶递推和双阶乘表达式见式~\eqref{eq:pei-wallis-integral}，混合正弦、余弦幂的公式见式~\eqref{eq:pei-mixed-wallis}。这些公式与上面的两个递推都依靠分部积分；遇到奇点时还须额外验证边界项的极限。

\begin{example}
计算 $\displaystyle J=\int_0^{\pi/2}\ln(\sin x)\,dx$，以及 $\displaystyle K=\int_0^\infty\frac{\ln x}{1+x^2}\,dx$。
\end{example}
\begin{solve}
因 $\sin x\sim x$，第一积分在零点附近与 $\ln x$ 相差有界量，故绝对收敛。由 $x\mapsto\pi/2-x$，$J=\int_0^{\pi/2}\ln(\cos x)\,dx$，所以
\[
 2J=\int_0^{\pi/2}\ln(\sin2x)\,dx-\frac\pi2\ln2
 =J-\frac\pi2\ln2,
 \qquad J=-\frac\pi2\ln2.
\]
第二积分在零点由 $|\ln x|$ 控制，在无穷远由 $(\ln x)/x^2$ 控制，故绝对收敛；换元 $x=1/t$ 得 $K=-K$，所以 $K=0$。这里先证收敛，才可将积分分拆、换元并移项，避免以形式对称掩盖发散。
\end{solve}

\subsection{比较判别与临界阶数}
\begin{theorem}
\label{thm:pei-int-improper-cauchy}
局部可积函数的反常积分 $\int_a^\infty f$ 收敛，当且仅当对每个 $\varepsilon>0$，存在 $A>a$，使对所有 $v>u>A$ 均有 $|\int_u^vf|<\varepsilon$。若 $f\ge0$，收敛又等价于截断积分 $\int_a^Rf$ 有上界。
\end{theorem}
\begin{proof}
第一条是实数极限存在的 Cauchy 准则用于函数 $R\mapsto\int_a^Rf$；第二条利用该函数随 $R$ 单调不减。有限反常端点的版本只需把远端尾区间改成趋向该端点的小区间。
\end{proof}
若尾部 $0\le f\le g$，则 $\int g$ 收敛推出 $\int f$ 收敛，$\int f$ 发散推出 $\int g$ 发散。若 $f/g\to c\in(0,\infty)$，二者同敛散；极限为零时只保留“小者可由大者收敛控制”的方向，极限为无穷时相反。反常点附近的等价无穷小可用于非负函数的敛散性判断，不能据此直接宣称两个收敛积分的值相同。

基本标尺是 $\int_1^\infty x^{-p}\,dx$ 当且仅当 $p>1$ 收敛，而 $\int_0^1x^{-p}\,dx$ 当且仅当 $p<1$ 收敛。临界幂次必须继续看对数：
\[
 \int_e^\infty\frac{dx}{x(\ln x)^q}<\infty\iff q>1,
 \qquad
 \int_0^{e^{-1}}\frac{dx}{x|\ln x|^q}<\infty\iff q>1.
\]
两式都由 $t=|\ln x|$ 化为无穷远的幂积分。若 $f>0$ 且 $\ln f(x)/\ln x\to-\alpha$，则 $\alpha>1$ 保证无穷远积分收敛，$\alpha<1$ 保证发散，$\alpha=1$ 不定；因为对任意小的 $\varepsilon>0$，最终有 $x^{-\alpha-\varepsilon}\le f(x)\le x^{-\alpha+\varepsilon}$。例如 $1/x$ 与 $1/(x\ln^2x)$ 都具有临界指数一，却有不同敛散性。

\begin{example}
讨论参数积分 $\displaystyle I_{\alpha,k}=\int_0^\pi\frac{(\sin x)^{\alpha-1}}{|1+k\cos x|^\alpha}\,dx$ 的收敛性，其中 $\alpha,k\in\mathbb R$。
\end{example}
\begin{solve}
若 $|k|<1$，分母处处远离零，两个端点均为 $x^{\alpha-1}$ 型，故当且仅当 $\alpha>0$ 收敛。若 $k=1$，零点附近要求 $\alpha>0$；但 $x\uparrow\pi$ 时 $1+\cos x\sim(\pi-x)^2/2$，被积函数为常数倍的 $(\pi-x)^{-\alpha-1}$，要求 $\alpha<0$，两者矛盾；$k=-1$ 由反射同样发散。若 $|k|>1$，存在唯一内部零点 $c$，且 $1+k\cos x$ 在 $c$ 处为一阶零点，故内部要求 $\alpha<1$，两端要求 $\alpha>0$。综上，收敛当且仅当 $|k|<1,\alpha>0$，或 $|k|>1,0<\alpha<1$。在各反常点分别比较，才能得到全部参数条件。
\end{solve}

\subsection{Dirichlet、Abel 判别与振荡抵消}
\begin{theorem}
\label{thm:pei-int-dirichlet-abel}
设 $f,g$ 在每个有限子区间上可积。Dirichlet 判别法：若 $F(R)=\int_a^Rf$ 有界，$g$ 单调且趋于零，则 $\int_a^\infty fg$ 收敛。Abel 判别法：若 $\int_a^\infty f$ 收敛，$g$ 单调且有界，则 $\int_a^\infty fg$ 收敛。
\end{theorem}
\begin{proof}
Dirichlet 情形可在必要时把 $g$ 换成 $-g$，设 $g\ge0$ 递减。若 $|F|\le C$，由第二积分中值定理，对 $B>A$ 有 $|\int_A^Bfg|\le2Cg(A)\to0$，从而满足 Cauchy 准则。Abel 情形中，$\int f$ 收敛保证所有足够远的 $|\int_A^tf|$ 统一小于 $\varepsilon$；对单调 $g$ 使用式~\eqref{eq:pei-int-bonnet}，得 $|\int_A^Bfg|\le2\sup|g|\varepsilon$，故也收敛。
\end{proof}
两种判别法利用振荡抵消，结论只保证收敛；要区分绝对与条件收敛，仍须另查绝对值积分。对周期振荡，$|\sin x|$ 与 $\sin^2x$ 每个周期都具有固定的正积分，故其乘以单调幂权重的积分可与相应幂积分比较；不能因 $\sin x$ 变号就把 $|\sin x|$ 的平均贡献也忽略。

\begin{example}
\label{ex:pei-int-oscillatory-power}
对实数 $\alpha,\beta$，分类 $\displaystyle\int_0^\infty x^\alpha\sin(x^\beta)\,dx$ 的敛散性。
\end{example}
\begin{solve}
若 $\beta=0$，被积函数为 $x^\alpha\sin1$，两个端点不能同时收敛。若 $\beta\ne0$，令 $t=x^\beta$，注意 $\beta<0$ 时积分方向反转，原积分化为
\[
 \frac1{|\beta|}\int_0^\infty t^{r-1}\sin t\,dt,
 \qquad r=\frac{\alpha+1}{\beta}.
\]
零点处与 $t^r$ 等价，要求 $r>-1$；无穷远处当 $r<1$ 时由 Dirichlet 判别收敛，当 $r\ge1$ 时取同号半周期，其积分绝对值不趋于零，违背 Cauchy 准则。绝对值积分在无穷远处收敛当且仅当 $r<0$：可逐周期用权重的上下界比较 $\int|\sin t|$ 与级数 $\sum n^{r-1}$。因此 $-1<r<0$ 时绝对收敛，$0\le r<1$ 时条件收敛，其余情况发散。
\end{solve}

\begin{example}
对 $p>0$，讨论 $\displaystyle\int_2^\infty\frac{\sin x}{x^p+\sin x}\,dx$。
\end{example}
\begin{solve}
分母在该区间严格为正，利用恒等式
\[
 \frac{\sin x}{x^p+\sin x}
 =\frac{\sin x}{x^p}-\frac{\sin^2x}{x^p(x^p+\sin x)}.
\]
第一项的积分对所有 $p>0$ 均由 Dirichlet 判别收敛；第二项非负，且与 $x^{-2p}\sin^2x$ 有固定比例的上下界，因此当且仅当 $2p>1$ 可积。原积分故当且仅当 $p>1/2$ 收敛。其绝对值与 $x^{-p}|\sin x|$ 可比较，所以当且仅当 $p>1$ 绝对收敛，$1/2<p\le1$ 时条件收敛。仅取分母的主项会漏掉不振荡的二阶贡献，正是本题的关键。
\end{solve}

\subsection{积分收敛与函数趋于零的关系}
函数趋于零不保证反常积分收敛，例如 $1/(1+x)$；反过来，收敛也不保证函数趋于零，甚至同时要求函数非负、连续也不够。具体地，在整数 $n\ge2$ 周围放置三角尖峰 $\psi_n(x)=\max(1-2^n|x-n|,0)$，令 $f=\sum_{n=2}^\infty\psi_n$。支集互不相交，和在每个有限区间上只有有限项，故 $f$ 连续非负；而 $\int_0^\infty f=\sum_{n=2}^\infty2^{-n}<\infty$，但 $f(n)=1$。这里尖峰越来越窄，说明必须补充控制局部振荡的条件。

\begin{theorem}
\label{thm:pei-int-integral-decay}
若 $f$ 在 $[a,\infty)$ 一致连续，且 $\int_a^\infty f$ 收敛，则 $f(x)\to0$。若 $f\ge0$ 单调不增且积分收敛，则更有 $xf(x)\to0$。
\end{theorem}
\begin{proof}
若第一条不成立，存在 $\varepsilon_0>0$ 及 $x_n\to\infty$ 使 $|f(x_n)|\ge\varepsilon_0$。一致连续性给出与 $n$ 无关的 $\delta>0$，使在 $[x_n,x_n+\delta]$ 上 $f$ 与 $f(x_n)$ 同号且绝对值至少为 $\varepsilon_0/2$，故该段积分的绝对值至少为 $\delta\varepsilon_0/2$，与积分的 Cauchy 准则矛盾。第二条由单调性得到
\[
 0\le\frac x2f(x)\le\int_{x/2}^xf(t)\,dt\longrightarrow0.
 \qedhere
\]
\end{proof}
导数有界保证一致连续，因而是第一条的一个充分条件。若 $f\in C^1[a,\infty)$，$\int f$、$\int f'$ 都收敛，则 $f(x)=f(a)+\int_a^xf'$ 先给出 $f$ 的有限极限；该极限若非零，$\int f$ 必发散，所以极限只能为零。若进一步已知 $xf(x)\ge0$ 单调不增，且 $\int f$ 收敛，则 $\int_{\sqrt x}^xf(t)\,dt\ge\tfrac12xf(x)\ln x$，故还有 $xf(x)\ln x\to0$。各个衰减结论依赖的单调对象不同，不宜混用。

\subsection{反常积分的极限：截断与尾部控制}
在无穷区间上，局部一致收敛仍不足以交换积分与极限：质量可能向无穷远移动。例如在 $[n,n+1]$ 上放一个面积为一的连续三角峰，其函数列在每个固定紧区间最终恒为零，但积分始终为一。一个直接可用的充分条件是：$f_n\to f$ 在每个 $[a,A]$ 上一致成立，且存在非负可积函数 $H$ 使 $|f_n|,|f|\le H$；将误差积分拆成 $[a,A]$ 和尾部，先使 $2\int_A^\infty H$ 小，再固定 $A$ 令 $n$ 大，即可得到 $\int f_n\to\int f$。这只是用截断估计证明的充分条件，不能省去尾部的一致控制。

\begin{proposition}
\label{prop:pei-int-abel-average}
若 $f$ 局部 Riemann 可积且 $f(x)\to L$（$x\to\infty$），则
\[
 \lim_{t\downarrow0}t\int_0^\infty e^{-tx}f(x)\,dx=L.
\]
若 $f,\phi$ 局部可积，$\phi>0$ 最终成立、$\int_a^x\phi\to\infty$，且 $f(x)/\phi(x)\to L$，则
\[
 \frac{\int_a^xf(t)\,dt}{\int_a^x\phi(t)\,dt}\longrightarrow L.
\]
\end{proposition}
\begin{proof}
第一式由 $t\int_0^\infty e^{-tx}\,dx=1$ 减去 $L$。选 $A$ 使 $x\ge A$ 时 $|f-L|<\varepsilon$，则尾部误差不超过 $\varepsilon$；固定段误差不超过 $t\int_0^A|f-L|\to0$。第二式对 $f-L\phi$ 同样分段，尾部绝对值不超过 $\varepsilon\int_A^x\phi$，固定段除以发散的分母趋于零。该证明只要求局部可积性，无需使用可能不适用的 L'Hospital 法则。
\end{proof}
取 $\phi\equiv1$ 得到积分的 Cesàro 平均 $x^{-1}\int_0^xf\to L$。对 $f\in C[0,1]$，类似的集中估计还给出
\[
 \lim_{x\downarrow0}x^\alpha\int_x^1\frac{f(t)}{t^{\alpha+1}}\,dt=\frac{f(0)}\alpha\quad(\alpha>0),
 \qquad
 \lim_{x\downarrow0}\frac{\int_x^1\frac{f(t)}t\,dt}{\ln(1/x)}=f(0).
\]
第一式减去 $f(0)x^\alpha\int_x^1t^{-\alpha-1}\,dt$，在 $[x,\delta]$ 上用连续性、在 $[\delta,1]$ 上用 $x^\alpha\to0$；第二式则把分母看成同一正权 $1/t$ 的积分。端点处的函数值决定主项，而区间远端只贡献低阶项。

\begin{proposition}
\label{prop:pei-int-gaussian-kernel}
设 $f\in C[a,b]$，$a<c<b$。则
\[
 \lim_{n\to\infty}\sqrt{\frac n\pi}\int_a^be^{-n(x-c)^2}f(x)\,dx=f(c).
\]
若 $c=a$ 或 $c=b$，极限为 $f(c)/2$。若 $f$ 在 $[a,\infty)$ 局部可积且 $\int_a^\infty|f|<\infty$，$g$ 为周期 $T$ 的 Riemann 可积函数，则
\[
 \lim_{\lambda\to\infty}\int_a^\infty f(x)g(\lambda x)\,dx
 =\frac1T\int_0^Tg\,\int_a^\infty f.
\]
\end{proposition}
\begin{proof}
Gaussian 核令 $u=\sqrt n(x-c)$ 后，积分范围向整条实轴扩张；在 $c$ 的小邻域内用 $f$ 的连续性，外部由 Gaussian 尾部控制，总质量趋于一。集中点在端点时仅保留一侧，质量为一半。第二式先在 $[a,A]$ 上应用定理~\ref{thm:pei-int-periodic-average}，尾部误差由 $2\sup|g|\int_A^\infty|f|$ 控制，先固定足够大的 $A$ 再令 $\lambda\to\infty$ 即得。
\end{proof}
例如周期函数 $\psi$ 的平均值为 $C$ 时，换元 $t=nu$ 得
\[
 n\int_n^\infty\frac{\psi(t)}{t^2}\,dt
 =\int_1^\infty\frac{\psi(nu)}{u^2}\,du\longrightarrow C.
\]
绝对可积的衰减权把无穷区间的周期振荡问题化回有限区间的平均定理。

\subsection{反常积分何时仍可用和式逼近}
\begin{proposition}
\label{prop:pei-int-singular-sums}
设 $f\ge0$ 在 $(0,1]$ 单调不增，且 $\int_0^1f<\infty$，则
\[
 \frac1n\sum_{k=1}^nf(k/n)\longrightarrow\int_0^1f(x)\,dx.
\]
对每个固定 $\theta\in(0,1)$，$n^{-1}\sum_{k=0}^{n-1}f((k+\theta)/n)$ 也有同一极限。
\end{proposition}
\begin{proof}
由单调性，$\int_{1/n}^1f\le n^{-1}\sum_{k=1}^nf(k/n)\le\int_0^1f$，两端趋于同一有限值。移位和不小于右端点和，且至多比它多 $f(\theta/n)/n$；又因 $f$ 递减，$0\le f(\theta/n)/n\le\theta^{-1}\int_0^{\theta/n}f\to0$，故移位和也收敛。
\end{proof}
若 $f$ 只在奇点附近单调，而远离奇点后正常可积，可分段应用上述结论与普通积分和极限。\emph[caution]{这不是任意反常可积函数都具有的 Riemann 和定义}；无界函数在极靠近奇点处可取任意大的标记值，任意标记和通常不能统一收敛。条件中有限积分也不可略去；和式有某种广义极限并不能反推通常的反常积分收敛。

\begin{example}
设 $0<a<d$，$A_n$、$G_n$ 分别为 $a,a+d,\ldots,a+(n-1)d$ 的算术平均与几何平均，求 $\lim G_n/A_n$。
\end{example}
\begin{solve}
记 $\theta=a/d\in(0,1)$。有 $A_n/(nd)=\theta/n+(n-1)/(2n)\to1/2$，且
\[
 \ln\frac{G_n}{nd}=\frac1n\sum_{k=0}^{n-1}\ln\frac{k+\theta}{n}
 \longrightarrow\int_0^1\ln x\,dx=-1.
\]
这里对 $-\ln x$ 应用命题~\ref{prop:pei-int-singular-sums}，不能把零点无界的对数直接当作闭区间上的 Riemann 可积函数。故 $G_n/(nd)\to e^{-1}$，最终 $G_n/A_n\to2/e$。
\end{solve}

\subsection{非负截断与几个综合推论}
若 $f\ge0$ 在 $(0,\infty)$ 连续，令 $f_M=\min(f,M)$，则 $\lim_{M\to\infty}\int_0^\infty f_M=\int_0^\infty f$，允许两边为 $+\infty$。在任意 $[\delta,R]$ 上，$M$ 足够大时 $f_M=f$，因而左端极限不小于每个截断积分；另一方面 $f_M\le f$，取所有截断区间的上确界即可证明等式。这一证明依赖非负性提供的单调结构，不能不加条件地用于变号函数或任意交换两个极限。

若连续正函数 $f$ 在 $[0,\infty)$ 上可积，Schwarz 不等式给出
\[
 R^2\le\left(\int_0^Rf\right)\left(\int_0^R\frac1f\right),
\]
所以 $\int_0^\infty1/f$ 必发散。这是把有限区间上的不等式用于无穷区间的一个例子。另一方面，仅有两个非负函数各自积分发散，并不能据比较判别推出其最小值积分也发散：例如 $f(x)=\max(\sin x,0)$、$g(x)=\max(-\sin x,0)$ 连续非负，两者的积分都发散，但 $\min(f,g)\equiv0$。比较判别所需的是逐点大小关系和正确的推理方向。

若 $f$ 在实轴上局部可积，$f(x)\to A$（$x\to+\infty$）、$f(x)\to B$（$x\to-\infty$），则对固定实数 $h$，通常的两侧反常积分满足
\[
 \int_{-\infty}^\infty\bigl(f(x+h)-f(x)\bigr)\,dx=h(A-B).
\]
在有限区间 $[u,v]$ 上，原积分等于 $\int_v^{v+h}f-\int_u^{u+h}f$；令两端分别趋于无穷即得结论，且分别固定另一端也证明两条尾积分各自收敛。因此这个恒等式不依赖对称主值，也不要求 $\int f$ 本身收敛。

% 教材来源：数学分析/textbooks/数分习题（分章原文）/数分习题-5.pdf。
% 按四节主题整合正文与代表性练习；非逐题转录。




